\chapter{Comparing maya to other TUI libraries}
\label{ch:comparison}

\begin{aside}
\maya{} is one of many TUI libraries. Each has its own
opinions; \maya{}'s opinions only make sense in context. This
chapter compares \maya{} to the major alternatives: ncurses,
notcurses, tui-rs / ratatui, Bubble Tea, Textual, FTXUI, Brick.
The point is not to declare a winner but to make the design
trade-offs explicit.
\end{aside}

\section{ncurses (C, 1980s)}

The grandparent of TUI libraries. Origin: BSD's curses, extended
by GNU. The model:

\begin{itemize}[leftmargin=*]
  \item Imperative: \code{move(y, x); addstr(s); refresh();}
  \item Window-oriented: every region is a \code{WINDOW*}.
  \item Manual layout: you compute coordinates.
  \item Manual input loop: \code{wgetch} in a \code{while}.
\end{itemize}

\textbf{Strengths.} Ubiquitous, stable, available everywhere. The
``every Unix machine has it'' guarantee.

\textbf{Weaknesses.} Imperative model fights composition. No
flexbox, no themes, no rich text. Truecolour is bolted on. Modern
keyboard handling (KKP, bracketed paste) requires manual escape
parsing.

\textbf{\maya{} contrast.} Declarative \code{view}, flexbox, type-
safe DSL. The trade: \maya{} needs a modern compiler; ncurses
runs on anything.

\section{notcurses (C, 2020)}

A modern reimagining of curses by Nick Black. Targets the
absolute highest fidelity of terminal output.

\textbf{Strengths.} Pixel-perfect rendering (cell-by-cell control).
First-class support for images via sixel and the Kitty graphics
protocol. Truecolour everywhere.

\textbf{Weaknesses.} Imperative API (similar to ncurses but
modernised). Steep learning curve. C API; the C++ wrapper is
thin.

\textbf{\maya{} contrast.} \maya{} is declarative and offers
fewer rendering capabilities (no images, no sixel). The trade:
\maya{}'s code is shorter and more testable; notcurses output is
visually richer.

\section{tui-rs / ratatui (Rust, 2019-)}

The Rust ecosystem's standard. Originally by Florian Dehau as
tui-rs; the community fork ratatui continues development.

\textbf{Strengths.} Declarative widget API. Strong Rust type
safety. Active community, dozens of community widgets.

\textbf{Weaknesses.} Rust-only. Widget composition uses
\code{Box<dyn Widget>} for dynamism; \maya{}'s sum-type model
avoids the indirection. Layout is constraint-list-based, not
flexbox-based, and feels less ergonomic for nested cases.

\textbf{\maya{} contrast.} Different languages, similar
declarative spirit. \maya{}'s DSL pushes more into compile time;
ratatui pushes more into runtime composition.

\section{Bubble Tea (Go, 2020)}

By Charm. The library that introduced the Elm architecture to
many programmers via Go.

\textbf{Strengths.} The Elm architecture made approachable.
Beautiful default styles via the Lip Gloss style library.
Excellent docs and examples.

\textbf{Weaknesses.} Go's lack of generics (until recently)
forced \code{interface\{\}} in many APIs. Rendering is line-based
(reasoned about as strings), not cell-based, so subtle Unicode
width bugs creep in.

\textbf{\maya{} contrast.} Same architectural family. \maya{}
is cell-based (graphemes + widths handled at the canvas level),
type-safer, and lower-level. Bubble Tea is faster to prototype in.

\section{Textual (Python, 2021)}

By the Rich author Will McGugan. Brings web-app ergonomics to
the terminal.

\textbf{Strengths.} CSS-like styling, including selectors and
themes. Reactive properties via Python descriptors. Beautiful
widget catalogue (Rich tables, syntax highlighting).

\textbf{Weaknesses.} Python overhead is real: a Textual app can
hit 5--10\% CPU at idle. The reactive system is implicit and
sometimes surprising.

\textbf{\maya{} contrast.} Different language ecosystems. Textual
appeals to web devs; \maya{} appeals to systems devs. Textual's
themes and styles are richer; \maya{}'s performance ceiling is
higher.

\section{FTXUI (C++, 2018-)}

The closest C++ analogue to \maya{}. By Arthur Sonzogni.

\textbf{Strengths.} Header-only. Functional-style API: components
are functions returning elements. Good widget catalogue.

\textbf{Weaknesses.} Uses inheritance for some component
hierarchies. Layout is less expressive than flexbox. The
component model is hookier than the Elm architecture (state lives
in closures, not in a plain model).

\textbf{\maya{} contrast.} Same language, different design
choices. FTXUI: closure-based state, classes for some
abstractions. \maya{}: value-based model, sum types throughout.
FTXUI is easier to learn; \maya{} scales further.

\section{Brick (Haskell, 2015)}

By Jonathan Daugherty. Haskell's Elm-inspired TUI library.

\textbf{Strengths.} Haskell's type system shines: the Elm
architecture is the natural shape. Pure functions everywhere.
Layout combinators are elegant.

\textbf{Weaknesses.} Haskell ecosystem; less mainstream than C++,
Go, or Python.

\textbf{\maya{} contrast.} Same architectural family. Brick
benefits from Haskell's type inference and pattern matching;
\maya{} works around C++'s limitations with concepts and visit.

\section{Where maya fits}

\maya{} occupies a specific niche:

\begin{itemize}[leftmargin=*]
  \item C++26 modern: NTTP-based DSL, concepts everywhere.
  \item Elm architecture: pure update, side effects as data.
  \item Cell-based renderer: SIMD diff, packed cells, integer
    layout.
  \item Compile-time guarantees: shape correctness checked at
    build, not runtime.
\end{itemize}

If you are a C++ shop that wants modern declarative TUIs with
strong static guarantees and high performance, \maya{} is in
your shortlist. If you need to deploy on legacy systems, want
to use Python or Go, or need image rendering, look at the
alternatives.

\section{Migration paths}

\subsection*{From ncurses}

The architectural shift is the hard part. Stop thinking in
``move cursor here, write text''; start thinking in ``what is
the model''. The rendering details (raw mode, escape sequences)
disappear into the runtime; you keep the business logic.

\subsection*{From Bubble Tea or Brick}

Mechanical: the model/message/update/view shape is the same.
You will rewrite Go's structs as C++ structs, channels as
\code{Cmd::task}, and the type system goes from inferred to
explicit. Most of your logic translates line by line.

\subsection*{From FTXUI}

Less mechanical: closures-as-state become explicit fields in
\code{Model}. The flow of data is more visible; the trade is
some up-front design work for testability later.

\subsection*{From Textual}

A bigger jump: dynamic Python to static C++. Worth it if you
need the performance; do not bother if Textual is fast enough.

\section{Choosing wisely}

A handful of questions to ask:

\begin{itemize}[leftmargin=*]
  \item \textbf{Language?} The library should match your
    ecosystem; you will not enjoy fighting bindings.
  \item \textbf{Performance?} If you need sub-millisecond
    frame budget, prefer C++ or Rust libraries. Python and
    Go have their own ceilings.
  \item \textbf{Type safety?} If you value compile-time
    catching of structural bugs, \maya{}, ratatui, and
    Brick are good fits.
  \item \textbf{Visual fidelity?} If you need images and
    fancy rendering, notcurses leads.
  \item \textbf{Familiarity?} If your team already knows one
    library, the switching cost may dominate the benefits.
\end{itemize}

There is no objectively best library. The best is the one
whose trade-offs match your project.

\section*{Workshop}

\begin{enumerate}[leftmargin=*]
  \item Pick two libraries from this chapter. Build a
    \emph{Hello World} in each. Compare the code.
  \item Identify three patterns from
    Chapter~\ref{ch:patterns} and find how each library
    expresses them.
  \item Choose a real project. List the requirements, then
    pick the library best suited to them. Defend the choice.
\end{enumerate}

\begin{recap}
\maya{} is one TUI library among many. ncurses is universal but
imperative; notcurses is high-fidelity; ratatui is the Rust
standard; Bubble Tea is the Elm story in Go; Textual is the
Python web-ish option; FTXUI is the C++ closure-based option;
Brick is the Haskell idiom. \maya{}'s position: modern C++,
Elm architecture, compile-time guarantees, high performance.
Pick by trade-off.
\end{recap}
